\paragraph{Tied rows and the main-block determinant}
Realize all the representations on lattices over one coefficient
discrete valuation ring \(\mathcal O\), taking the direct sum of the
main-block and character lattices for the genuine lifts.
Indeed, over \(\overline{\mathbb Q}_\ell\) form the candidate representation in the
product of the main matrix algebra and the \(t'\) scalar algebras. The trace on its
diagonal realization equals the trace of \(N_x\). Its span is semiprime.
If \(yU(g)y=0\) for every \(g\), then on the main summand this holds for every matrix in
place of \(U(g)\), since the main lattice spans the full integral matrix algebra.
Taking the matrix unit \(E_{ba}\) forces \(y_{ab}^2=0\).
On each scalar block, take \(g=1\) to obtain the same conclusion.
The trace form on the span of \(N_x\) is nondegenerate by the attached representation theorem.
Brauer--Nesbitt over the algebraically closed field of characteristic zero
therefore gives an intertwining change of basis with the diagonal realization.
In detail, the combined image spans a semisimple algebra by the same semiprime test
on the two factors. In its Wedderburn decomposition, equal traces give equal ranks of
the images of a rank-one idempotent from each simple factor, so the multiplicities of
the simple modules agree. Conjugating to this realization gives the required integral lattice.
\leanlinks{\href{\upstream/grh/src/New/Congruences/FirstOrder/TraceRelations.lean\#L402}{The direct-sum lattice from equality of traces}}

Write \(\rho\) for the twisted main block at \(\pi\),
\(\gamma=\gamma_\pi\), and \(a_{\mathrm c}=a_{\mathrm{cyc}}^{-1}\).
Use \(v_{\mathcal O}\) with a uniformizer of valuation \(1\).
Dividing lengths in this convention by \(v_{\mathcal O}(\ell)\)
gives the normalization corresponding to \(v_\ell\).
Set \(\mathcal S=\mathcal O[(X_j)_{j\in J}]\) and
\(\mathcal A=\mathcal S[G_E]\), identifying \(g\in G_E\) with its
basis element in \(\mathcal A\).

The \(\mathcal O\)-span of the families \((N_x(g)-1)_x\) is finitely
generated. Choose generators indexed by a finite family \((g_i)_{i\in I}\),
including the old inertial elements and the detection points used below.
For each \(g\) choose coordinates \(z(g)=(z_i(g))_{i\in I}\in\mathcal O^I\)
satisfying
\[
 (N_x(g)-1)_x=\sum_{i\in I} z_i(g)(N_x(g_i)-1)_x .
\]
For a row \(b\in\mathcal A^I\) write \(b\cdot\mathbf y=\sum_i b_i(g_i-1)\).
Let \(\mathcal I_\mathrm{aug}\) be the augmentation ideal of \(\mathcal A\), and put
\[
 c_\gamma=(\zeta_\pi-1)^{-1},\qquad n_\gamma=c_\gamma(\gamma-1),
 \qquad \rho_0=n_\gamma+(1-n_\gamma)(a_{\mathrm c}-1).
\]
The main block \(\rho\) sends \(\rho_0\) to the identity, and the
character \(\chi_a\) sends it to \(\chi_a(a_{\mathrm c})-1\).

Let \(N_\mathrm{row}\) be the module of rows for which \(b\cdot\mathbf y\)
belongs to the left ideal generated by four sets. They are the
common kernel of the \(N_x\), \(p_\ast\mathcal I_\mathrm{aug}\),
the new-prime ties \(\tau-1-X_jH_j(\phi)(\phi-1)\) and their
counterparts at the inverse conjugates, and
\(\mathcal I_\mathrm{aug}\rho_0\).
The row \(z(g)\) has the same value as \(g-1\) under every \(N_x\).
Hence \(N_\mathrm{row}\) contains
\[
 \begin{gathered}
 z(gh)-z(g)-g z(h),\qquad z(g_i)-\mathbf e_i,\qquad
 \sum_g\theta_g z(g)\quad
     \left(\sum_g\theta_g (N_x(g)-1)=0\text{ for all }x\right),\\
 p_\ast b\quad (b\in\mathcal A^I),\qquad
 z(\tau)-X_jH_j(\phi)z(\phi).
 \end{gathered}
\]
The first three rows give zero under every \(N_x\). The last row
evaluated on \(\mathbf y\) gives the tie itself under every \(N_x\).
The row \(u_\gamma=c_\gamma z(\gamma)+(1-n_\gamma)z(a_{\mathrm c})\)
evaluated on \(\mathbf y\) gives \(\rho_0\) under every \(N_x\).
Thus \((g-1)u_\gamma\in N_\mathrm{row}\) for every \(g\).
\leanlinks{\href{\upstream/grh/src/New/Congruences/FirstOrder/RibetMatrices.lean\#L175}{Tied rows and the relations fixing coboundaries}}

Let \(\varrho_0\colon\mathcal A\to M_{r'}(\mathcal O)\) be evaluation
on the main block at \(\pi\) and at \(X=0\). It is surjective since the
image of \(\rho\) spans the full matrix algebra. Put
\(K_{\mathcal O}=\operatorname{Frac}(\mathcal O)\).
Suppose the vectors \(w_i\in(K_{\mathcal O}/\mathcal O)^{r'}\)
satisfy \(\sum_i\varrho_0(b_i)w_i=0\) for every \(b\in N_\mathrm{row}\),
and set \(z_w(g)=\sum_i z_i(g)w_i\).
The first three kinds of row give the cocycle rule for \(\rho\),
the equality \(z_w(g_i)=w_i\), and the linear relations among the
families \(N_x(g)-1\). The tie at \(X=0\) gives \(z_w(\tau)=0\).
The row \((\gamma-1)u_\gamma\) has image \(z(\gamma)\) because
\(\rho(\gamma)-1=(\zeta_\pi-1)I\) and \(1-n_\gamma\) acts by zero.
Thus \(z_w(\gamma)=0\). The assignment \(w\mapsto z_w\) embeds
the annihilator of the rows into this divisible-coefficient cocycle
module. Write \(b_{\mathrm{main}}\) for the normalized length bound on that
module obtained by the local comparison below.
\leanlinks{\href{\upstream/grh/src/New/Congruences/FirstOrder/FittingBlock.lean\#L619}{From the annihilator of the rows to cocycles}
  \leanlinksep \href{\upstream/grh/src/New/Congruences/FirstOrder/RibetMatrices.lean\#L539}{Specialization on the main block}}

To pass from the length of the annihilator to a determinant,
choose \(p_{\mathrm{row}}\in\mathcal O\) with \(v_\ell(p_{\mathrm{row}})>b_{\mathrm{main}}\), for example a
power of the uniformizer of exponent greater than \(v_{\mathcal O}(\ell)b_{\mathrm{main}}\).
Set \(d=\#I\cdot r'\), and let \(\sigma(b)\in\mathcal O^d\)
be the sum of the \(r'\) rows of the block row \((\varrho_0(b_i))_{i\in I}\).
The module \(\sigma(N_\mathrm{row})\) contains each individual row.
Indeed, multiplying on the left by a preimage of the matrix unit
\(E_{cc}\) makes \(\sigma\) select row \(c\). Form the full-rank row lattice
\[
 L=p_{\mathrm{row}}\mathcal O^d+\sigma(N_\mathrm{row}).
\]
It contains \(p_{\mathrm{row}}\mathcal O^d\), so Smith normal form identifies the
length of its annihilator in \((\mathcal O/p_{\mathrm{row}})^d\) with
\(\operatorname{length}(\mathcal O^d/L)\). Componentwise division
by \(p_{\mathrm{row}}\) embeds that annihilator into the one with coefficients
in \(K_{\mathcal O}/\mathcal O\). The matrix \(\mathcal M\) whose
rows are a basis of \(L\) therefore satisfies
\(v_\ell(\det\mathcal M)\le b_{\mathrm{main}}\). Replacing each row by a row
\(\sigma(b)\) gives a congruent matrix modulo \(p_{\mathrm{row}}\), and its
determinant has the same valuation since \(v_\ell(p_{\mathrm{row}})>b_{\mathrm{main}}\).

Group the rows in blocks of \(r'\). For each \(c\), let \(F_c\) be
the square matrix with ones in row \(c\) and zero elsewhere.
Multiply each of the corresponding \(b\)'s by a preimage of \(F_c\)
and sum over \(c\). The resulting block row has the designated
\(\sigma(b)\) in each scalar row. This gives a square matrix \(B_0\)
indexed by \(I\) with rows in \(N_\mathrm{row}\). Writing \(D_x(X)\)
for its determinant on the main block, we have
\begingroup\postdisplaypenalty=10000
\[
 v_\ell(D_\pi(0))=v_\ell(\det\mathcal M)\le b_{\mathrm{main}} .
\]
\leanlinks{\href{\upstream/grh/src/Lib/General/Algebra/BlockFittingBound.lean\#L31}{The DVR row lattice and lifting its block rows}
  \leanlinksep \href{\upstream/grh/src/Lib/General/Algebra/BlockFittingBound.lean\#L152}{Passage from divisible coefficients to a finite quotient}}
\endgroup
